\section{智能和产品都重要}

本章讨论第二关键词：产品。过去三年行业一直极度上头地探索智能上限，但现在产品重新变得重要。ChatGPT 身上有很多非技术性壁垒，而 Coding 工具或模型公司如果只有技术壁垒，很容易被追赶。OpenAI 的强处在于平衡：一边探索智能上限，一边把智能红利转化成产品流量和品牌心智。

\begin{figure}[H]
\centering
\includegraphics[width=0.96\textwidth]{intelligence-product-balance.png}
\caption{智能和产品都重要：智能上限探索必须转成流量、品牌和用户心智。自制概念图，依据 00:33:35--00:38:52 对谈内容整理。}
\end{figure}

\begin{knowledgebox}{读图：智能不是自动等于产品}
模型能力提供可能性，产品体验把可能性变成默认使用，流量和品牌心智让用户持续回来，商业化把智能变成公司能力。缺任何一环，都可能只是技术 demo。
\end{knowledgebox}

\subsection{模型壁垒与产品壁垒}

上一节强调产品重新变重要，本节拆开两类壁垒。模型壁垒来自 benchmark、coding、agentic、上下文和推理能力；产品壁垒来自用户习惯、品牌、入口、数据回流和商业化。模型壁垒会被追赶，产品壁垒也会被吞并，但二者叠加才形成强公司。

\begin{figure}[H]
\centering
\includegraphics[width=0.94\textwidth]{model-vs-product-moat.png}
\caption{模型壁垒 vs 产品壁垒：技术壁垒会被追赶，产品壁垒包含流量、品牌和默认入口。自制概念图，依据 00:33:35--00:38:52 对谈内容整理。}
\end{figure}

\begin{warningbox}{不要只看技术壁垒}
AI 产品的技术壁垒可能在 3 到 6 个月内被追上。如果没有数据回流、用户习惯、渠道、品牌或工作流沉淀，技术领先很容易变成短暂窗口。
\end{warningbox}

\begin{knowledgebox}{产品壁垒的非技术部分}
\begin{tabularx}{\textwidth}{p{0.22\textwidth}p{0.34\textwidth}X}
\toprule
壁垒 & 表现 & 为什么难复制 \\
\midrule
默认入口 & 用户每天自然打开 & 需要长期品牌和使用习惯。 \\
工作流嵌入 & 产品进入团队流程 & 替换成本高，数据持续回流。 \\
品牌心智 & 用户把某类任务默认交给它 & 首个 Magic Moment 会被长期记住。 \\
商业化 & 有稳定付费和渠道 & 不只是模型强，还要懂销售和包装。 \\
\bottomrule
\end{tabularx}
\end{knowledgebox}

\subsection{本章小结}

智能和产品的关系变得更平衡。模型公司需要产品化，产品公司也需要模型理解；单纯技术或单纯产品都不够。

